% chapters/14/section-05.tex
\section{Program Animasi dan Antarmuka}
\label{sec:program-aset}

Dua skrip berikut melengkapi level pada Bab~XIII: satu menghubungkan keadaan
pemain dengan animasinya, dan satu menampilkan skor serta nyawa.

\subsection{Skrip Animasi}

\kodecsharp{unity/ch14_kontrol_animasi.cs}{Menghubungkan keadaan pemain dengan animasinya}{lst:ch14-animasi}

Skrip itu tidak memutar animasi apa pun. Ia hanya menetapkan
nilai parameter, sedangkan peralihan animasinya dikerjakan mesin menurut
transisi yang Anda gambar pada jendela \istilah{Animator}. Pembagian pekerjaan
itu sama dengan pembagian pada Bab~IX: kode menentukan keadaan, sedangkan
peralihannya ditentukan oleh aturan yang telah Anda susun.

Perhatikan pula cara arah hadap gambar dibalik, yaitu dengan mengubah
\code{transform.localScale}. Cara itu dipakai alih-alih memutar bendanya,
sebab memutar gambar dua dimensi sebesar 180 derajat akan membuatnya terbalik
atas-bawah.

\subsection{Skrip Antarmuka}

\kodecsharp{unity/ch14_tampilan_hud.cs}{Menampilkan skor dan nyawa pada antarmuka}{lst:14}

Skrip itu memuat dua hal yang telah Anda kenal: fungsi publik yang dipanggil
dari skrip lain, sebagaimana fungsi \code{TambahSkor}, dan pemeriksaan rujukan
yang belum dihubungkan.

Perhatikan pemakaian \code{enabled = false} ketika rujukan antarmukanya belum
dihubungkan. Cara itu mematikan skripnya, sehingga kesalahan yang sama tidak
terulang pada setiap bingkai. Kebiasaan menghentikan pekerjaan yang tidak
dapat dilanjutkan lebih baik daripada membiarkannya berjalan dan
menghasilkan kesalahan berulang.

\subsection{Percobaan yang Disarankan}

\begin{table}[htbp]
\centering
\caption{Percobaan pada animasi dan antarmuka}
\label{tab:percobaan-ch14}
{\setlength{\tabcolsep}{4pt}\footnotesize%
\begin{tabularx}{\textwidth}{@{}L{5.4cm} Y@{}}
\toprule
\headrow \textbf{Yang diubah dan caranya} & \textbf{Yang sebaiknya Anda amati} \\
\midrule
\istilah{Pixels Per Unit} pada aset pemain, diubah dari 100 menjadi 50. &
Gambar pemain menjadi dua kali lebih besar pada scene, sedangkan bentuk
batasnya tidak ikut berubah. Sesuaikan bentuk batasnya sesudah percobaan. \\
\addlinespace
Nama parameter pada kode, diubah sehingga tidak sama dengan yang ada pada
jendela \istilah{Animator}. &
Animasi tidak pernah berganti, dan program tetap berjalan tanpa satu pun pesan
kesalahan. Kembalikan sesudah percobaan. \\
\addlinespace
Nilai ambang kecepatan pada transisi, dinaikkan menjadi 5. &
Animasi berlari baru muncul ketika pemain benar-benar cepat. Perhatikan
perbedaan rasa geraknya. \\
\addlinespace
Rujukan teks pada antarmuka, diputus dari \istilah{Inspector}. &
Pesan peringatan muncul pada panel \istilah{Console}, dan antarmukanya tidak
diperbarui. Kembalikan sesudah percobaan. \\
\bottomrule
\end{tabularx}}
\end{table}

\begin{catatabox}
Percobaan kedua adalah yang paling penting untuk dicoba, sebab kesalahannya
\textbf{tidak menimbulkan pesan apa pun}. Kebiasaan memakai peubah tetap untuk
nama parameter, sebagaimana pada skrip di atas, mencegah kesalahan itu sejak
awal.
\end{catatabox}
